\section*{Registro y protocolo}

El protocolo se registró en OSF (\url{https://osf.io/5usf3/}; DOI 10.17605/OSF.IO/5USF3)~\cite{OSF2026Protocolo} antes de la búsqueda oficial, y está disponible en el registro. Las enmiendas A1--A6 posteriores al registro se describen en la \cref{tab:enmiendas}. Las enmiendas A1--A3 ya constan en OSF; las enmiendas A4--A6 se actualizarán en el registro antes de la difusión del manuscrito o de su envío a una revista.

\section*{Financiamiento}

\RSLFinanciamiento

\section*{Conflictos de interés}

La autora desarrolla su proyecto de tesis de pregrado, para optar el título profesional, sobre un tema relacionado, por lo que puede tener interés en que exista una brecha de investigación. Para limitar ese sesgo, el criterio de brecha y la regla del indicador D4 se registraron antes de examinar los datos.

\section*{Disponibilidad de datos y código}

El repositorio público del proyecto (\url{\RSLRepositorio}) contiene el protocolo, las cadenas y el registro de búsquedas, las exportaciones de arXiv y ACL Anthology, los registros de cribado con sus decisiones y motivos, los archivos de extracción y de calidad, los \en{scripts} de síntesis, las figuras y las fuentes de este artículo, con código bajo licencia MIT y datos bajo \RSLLicenciaDatos. Los textos completos de los estudios no se comparten por razones de derechos de autor. Las exportaciones en bruto de Scopus, IEEE Xplore, ACM DL y Semantic Scholar, cuyos términos de uso restringen la redistribución, y los resúmenes de los registros se omiten; un manifiesto lista cada archivo omitido con su suma \RSLSuma.

\section*{Declaración de uso de IA}

Se emplearon modelos de inteligencia artificial para el cribado, la extracción de datos y la redacción, bajo la verificación de la autora, según se describe en la \cref{sec:met-seleccion}. \pendiente{declaración completa con /ars-disclosure (fase 13)}
